% GENERATED FILE. Edit canonical lessons, then run npm run generate:books.
% canonical-source-hash: fnv1a64:ba08419173c381e1
% canonical-lessons: KA-C375-sibira, KA-C375-nirasrita, KA-C375-santapa, KA-C375-hastantarisu, KA-C375-dhavisu, KA-R375-review-369-375, KA-R375-review-final-375

\chapter{The Relief Camp}
\label{ch:ka-the-relief-camp}
\hlchaptermodality{375}

\begin{chapteropening}
\textbf{By the end of this chapter:} \emph{I can say a camp, a refugee, condolence, to hand over and to rush.}

Complete the last lesson of chapter 375: I can say a camp, a refugee, condolence, to hand over and to rush.
\end{chapteropening}

\section[\texorpdfstring{\kn{ಶಿಬಿರ} (śibira)}{ಶಿಬಿರ (śibira)}]{\kn{ಶಿಬಿರ} (śibira) — a camp}

\label{lesson:KA-C375-sibira}



\begin{quote}
Before the new one: say the Kannada for a noise, then the Kannada for to kill.
\end{quote}

\subsection*{You'll want to know: \kn{ಶಿಬಿರ}}
\textbf{\kn{ಶಿಬಿರ}} — \emph{śibira} — \textquotedblleft{}a camp\textquotedblright{}.

\begin{quote}
\kn{ಪ್ರವಾಹದ} \kn{ನಂತರ} \kn{ಶಾಲೆಯಲ್ಲಿ} \kn{ಒಂದು} \kn{ಶಿಬಿರ} \kn{ತೆರೆದರು}.
\end{quote}

\begin{quote}
\emph{pravāhada nantara śāleyalli ondu śibira teredaru.}
\end{quote}

\begin{quote}
After the flood they opened a camp in the school.
\end{quote}

Plural: \kn{ಶಿಬಿರಗಳು}, \emph{śibiragaḷu}.

Also a training or health camp: \kn{ಆರೋಗ್ಯ} \kn{ಶಿಬಿರ}.

\begin{cousinweb}
\textbf{\kn{ಶಿಬಿರ}} (\emph{śibira}) is borrowed from Sanskrit \emph{śibira}, \textquotedblleft{}a royal camp, a tent, an entrenchment\textquotedblright{} (Monier-Williams).
\end{cousinweb}

\begin{grammarlens}[title={What you've built}]
One more word for telling what happened.
\end{grammarlens}

\subsection*{Your turn}
\begin{itemize}
\raggedright
  \item \emph{Say it:} \emph{śibira}
  \item \emph{Say it:} \emph{śibira}, once more
  \item \emph{Say it:} \emph{pravāhada nantara śāleyalli ondu śibira teredaru.} — After the flood they opened a camp in the school.
  \item \emph{Recall:} say the Kannada for a noise, then the Kannada for to kill, then say \emph{śibira} again
\end{itemize}

\begin{tcolorbox}[breakable,colback=teal!4,colframe=teal!35!black,title={Before you move on}]
What does \kn{ಶಿಬಿರ} mean? (\textquotedblleft{}A camp\textquotedblright{}.) Say it once more.
\end{tcolorbox}
\section[\texorpdfstring{\kn{ನಿರಾಶ್ರಿತ} (nirāśrita)}{ನಿರಾಶ್ರಿತ (nirāśrita)}]{\kn{ನಿರಾಶ್ರಿತ} (nirāśrita) — a refugee; homeless}

\label{lesson:KA-C375-nirasrita}



\begin{quote}
Before the new one: say the Kannada for to kill, then the Kannada for a camp.
\end{quote}

\subsection*{You'll want to know: \kn{ನಿರಾಶ್ರಿತ}}
\textbf{\kn{ನಿರಾಶ್ರಿತ}} — \emph{nirāśrita} — \textquotedblleft{}a refugee; homeless\textquotedblright{}.

\begin{quote}
\kn{ಶಿಬಿರದಲ್ಲಿ} \kn{ಐನೂರು} \kn{ನಿರಾಶ್ರಿತರು} \kn{ಇದ್ದರು}.
\end{quote}

\begin{quote}
\emph{śibiradalli ainūru nirāśritaru iddaru.}
\end{quote}

\begin{quote}
There were five hundred refugees in the camp.
\end{quote}

Plural: \kn{ನಿರಾಶ್ರಿತರು}, \emph{nirāśritaru}.

\kn{ನಿರ್}-, without, and \kn{ಆಶ್ರಿತ}, sheltered: one without shelter.

\begin{cousinweb}
\textbf{\kn{ನಿರಾಶ್ರಿತ}} (\emph{nirāśrita}) is borrowed from Sanskrit: \emph{nir-}, \textquotedblleft{}without\textquotedblright{}, with \emph{āśrita}, \textquotedblleft{}sheltered, depending on\textquotedblright{} (Monier-Williams).
\end{cousinweb}

\begin{grammarlens}[title={What you've built}]
One more word for telling what happened.
\end{grammarlens}

\subsection*{Your turn}
\begin{itemize}
\raggedright
  \item \emph{Say it:} \emph{nirāśrita}
  \item \emph{Say it:} \emph{nirāśrita}, once more
  \item \emph{Say it:} \emph{śibiradalli ainūru nirāśritaru iddaru.} — There were five hundred refugees in the camp.
  \item \emph{Recall:} say the Kannada for to kill, then the Kannada for a camp, then say \emph{nirāśrita} again
\end{itemize}

\begin{tcolorbox}[breakable,colback=teal!4,colframe=teal!35!black,title={Before you move on}]
What does \kn{ನಿರಾಶ್ರಿತ} mean? (\textquotedblleft{}A refugee; homeless\textquotedblright{}.) Say it once more.
\end{tcolorbox}
\section[\texorpdfstring{\kn{ಸಂತಾಪ} (santāpa)}{ಸಂತಾಪ (santāpa)}]{\kn{ಸಂತಾಪ} (santāpa) — condolence; grief}

\label{lesson:KA-C375-santapa}



\begin{quote}
Before the new one: say the Kannada for a camp, then the Kannada for a refugee.
\end{quote}

\subsection*{You'll want to know: \kn{ಸಂತಾಪ}}
\textbf{\kn{ಸಂತಾಪ}} — \emph{santāpa} — \textquotedblleft{}condolence; grief\textquotedblright{}.

\begin{quote}
\kn{ಊರಿನವರು} \kn{ಸಂತ್ರಸ್ತರ} \kn{ಕುಟುಂಬಗಳಿಗೆ} \kn{ಸಂತಾಪ} \kn{ತಿಳಿಸಿದರು}.
\end{quote}

\begin{quote}
\emph{ūrinavaru santrastara kuṭumbagaḷige santāpa tiḷisidaru.}
\end{quote}

\begin{quote}
The townspeople offered condolences to the victims' families.
\end{quote}

A condolence meeting is a \kn{ಸಂತಾಪ} \kn{ಸಭೆ}.

\begin{cousinweb}
\textbf{\kn{ಸಂತಾಪ}} (\emph{santāpa}) is borrowed from Sanskrit \emph{saṃtāpa}, \textquotedblleft{}great heat\textquotedblright{}, and \textquotedblleft{}affliction, pain, sorrow\textquotedblright{} (Monier-Williams).
\end{cousinweb}

\begin{grammarlens}[title={What you've built}]
One more word for telling what happened.
\end{grammarlens}

\subsection*{Your turn}
\begin{itemize}
\raggedright
  \item \emph{Say it:} \emph{santāpa}
  \item \emph{Say it:} \emph{santāpa}, once more
  \item \emph{Say it:} \emph{ūrinavaru santrastara kuṭumbagaḷige santāpa tiḷisidaru.} — The townspeople offered condolences to the victims' families.
  \item \emph{Recall:} say the Kannada for a camp, then the Kannada for a refugee, then say \emph{santāpa} again
\end{itemize}

\begin{tcolorbox}[breakable,colback=teal!4,colframe=teal!35!black,title={Before you move on}]
What does \kn{ಸಂತಾಪ} mean? (\textquotedblleft{}Condolence; grief\textquotedblright{}.) Say it once more.
\end{tcolorbox}
\section[\texorpdfstring{\kn{ಹಸ್ತಾಂತರಿಸು} (hastāntarisu)}{ಹಸ್ತಾಂತರಿಸು (hastāntarisu)}]{\kn{ಹಸ್ತಾಂತರಿಸು} (hastāntarisu) — to hand over}

\label{lesson:KA-C375-hastantarisu}



\begin{quote}
Before the new one: say the Kannada for a refugee, then the Kannada for condolence.
\end{quote}

\subsection*{You'll want to know: \kn{ಹಸ್ತಾಂತರಿಸು}}
\textbf{\kn{ಹಸ್ತಾಂತರಿಸು}} — \emph{hastāntarisu} — \textquotedblleft{}to hand over\textquotedblright{}.

\begin{quote}
\kn{ಸ್ವಯಂಸೇವಕರು} \kn{ಊಟದ} \kn{ಪೊಟ್ಟಣಗಳನ್ನು} \kn{ಅಧಿಕಾರಿಗಳಿಗೆ} \kn{ಹಸ್ತಾಂತರಿಸಿದರು}.
\end{quote}

\begin{quote}
\emph{svayaṃsēvakaru ūṭada poṭṭaṇagaḷannu adhikārigaḷige hastāntarisidaru.}
\end{quote}

\begin{quote}
The volunteers handed over the food parcels to the officials.
\end{quote}

Past (I …): \kn{ಹಸ್ತಾಂತರಿಸಿದೆ}, \emph{hastāntariside}.

What is handed over takes -\kn{ಅನ್ನು}; who receives it takes -\kn{ಗೆ}.

\begin{cousinweb}
\textbf{\kn{ಹಸ್ತಾಂತರಿಸು}} (\emph{hastāntarisu}) is built from Sanskrit \emph{hasta}, \textquotedblleft{}the hand\textquotedblright{} (CDIAL 14024), and \emph{antara}, \textquotedblleft{}the inside, another\textquotedblright{} (CDIAL 357): into another's hand, with the Kannada verb ending -\emph{isu}.
\end{cousinweb}

\begin{grammarlens}[title={What you've built}]
One more word for telling what happened.
\end{grammarlens}

\subsection*{Your turn}
\begin{itemize}
\raggedright
  \item \emph{Say it:} \emph{hastāntarisu}
  \item \emph{Say it:} \emph{hastāntarisu}, once more
  \item \emph{Say it:} \emph{svayaṃsēvakaru ūṭada poṭṭaṇagaḷannu adhikārigaḷige hastāntarisidaru.} — The volunteers handed over the food parcels to the officials.
  \item \emph{Recall:} say the Kannada for a refugee, then the Kannada for condolence, then say \emph{hastāntarisu} again
\end{itemize}

\begin{tcolorbox}[breakable,colback=teal!4,colframe=teal!35!black,title={Before you move on}]
What does \kn{ಹಸ್ತಾಂತರಿಸು} mean? (\textquotedblleft{}To hand over\textquotedblright{}.) Say it once more.
\end{tcolorbox}
\section[\texorpdfstring{\kn{ಧಾವಿಸು} (dhāvisu)}{ಧಾವಿಸು (dhāvisu)}]{\kn{ಧಾವಿಸು} (dhāvisu) — to rush; to hasten}

\label{lesson:KA-C375-dhavisu}



\begin{quote}
Before the new one: say the Kannada for condolence, then the Kannada for to hand over.
\end{quote}

\subsection*{You'll want to know: \kn{ಧಾವಿಸು}}
\textbf{\kn{ಧಾವಿಸು}} — \emph{dhāvisu} — \textquotedblleft{}to rush; to hasten\textquotedblright{}.

\begin{quote}
\kn{ಸುದ್ದಿ} \kn{ಕೇಳಿ} \kn{ವೈದ್ಯರು} \kn{ಶಿಬಿರಕ್ಕೆ} \kn{ಧಾವಿಸಿದರು}.
\end{quote}

\begin{quote}
\emph{suddi kēḷi vaidyaru śibirakke dhāvisidaru.}
\end{quote}

\begin{quote}
Hearing the news, the doctors rushed to the camp.
\end{quote}

Past (I …): \kn{ಧಾವಿಸಿದೆ}, \emph{dhāviside}.

Where one rushes takes -\kn{ಗೆ}: \kn{ನೆರವಿಗೆ} \kn{ಧಾವಿಸು}, to rush to help.

\begin{cousinweb}
\textbf{\kn{ಧಾವಿಸು}} (\emph{dhāvisu}) is built on Sanskrit \emph{dhāv}, \textquotedblleft{}to run, to flow\textquotedblright{} (Monier-Williams), with the Kannada verb ending -\emph{isu}.
\end{cousinweb}

\begin{grammarlens}[title={What you've built}]
One more word for telling what happened.
\end{grammarlens}

\subsection*{Your turn}
\begin{itemize}
\raggedright
  \item \emph{Say it:} \emph{dhāvisu}
  \item \emph{Say it:} \emph{dhāvisu}, once more
  \item \emph{Say it:} \emph{suddi kēḷi vaidyaru śibirakke dhāvisidaru.} — Hearing the news, the doctors rushed to the camp.
  \item \emph{Recall:} say the Kannada for condolence, then the Kannada for to hand over, then say \emph{dhāvisu} again
\end{itemize}

\begin{tcolorbox}[breakable,colback=teal!4,colframe=teal!35!black,title={Before you move on}]
What does \kn{ಧಾವಿಸು} mean? (\textquotedblleft{}To rush; to hasten\textquotedblright{}.) Say it once more.
\end{tcolorbox}
\section[(dialogue)]{Telling what happened, the last stretch}

\label{lesson:KA-R375-review-369-375}



\begin{quote}
Say the Kannada for at last and with one's own eyes.
\end{quote}

\subsection*{Your turn}
Say these aloud:

\begin{itemize}
\raggedright
  \item at last, with one's own eyes, unexpected, a jail, to appear
  \item journalism, a documentary, a reader, a listener, notorious
  \item a wedding pendant, a palanquin, a chieftain, a mansion, to set foot
  \item exile in the forest, to impose, banishment, former, theft
  \item a rise, a downfall, the past, a narrative, a triumph
  \item magic, a miracle, false, a noise, to kill
  \item a camp, a refugee, condolence, to hand over, to rush
\end{itemize}

\begin{tcolorbox}[breakable,colback=teal!4,colframe=teal!35!black,title={Before you move on}]
At last? (\textbf{\kn{ಅಂತೂ}}, \emph{antū}.) With one's own eyes? (\textbf{\kn{ಕಣ್ಣಾರೆ}}, \emph{kaṇṇāre}.) Unexpected? (\textbf{\kn{ಅನಿರೀಕ್ಷಿತ}}, \emph{anirīkṣita}.) A jail? (\textbf{\kn{ಜೈಲು}}, \emph{jailu}.) To appear? (\textbf{\kn{ಕಾಣಿಸಿಕೊ}}, \emph{kāṇisiko}.) Journalism? (\textbf{\kn{ಪತ್ರಿಕೋದ್ಯಮ}}, \emph{patrikōdyama}.) A documentary? (\textbf{\kn{ಸಾಕ್ಷ್ಯಚಿತ್ರ}}, \emph{sākṣyacitra}.) A reader? (\textbf{\kn{ಓದುಗ}}, \emph{ōduga}.) A listener? (\textbf{\kn{ಕೇಳುಗ}}, \emph{kēḷuga}.) Notorious? (\textbf{\kn{ಕುಖ್ಯಾತ}}, \emph{kukhyāta}.) A wedding pendant? (\textbf{\kn{ತಾಳಿ}}, \emph{tāḷi}.) A palanquin? (\textbf{\kn{ಪಲ್ಲಕ್ಕಿ}}, \emph{pallakki}.) A chieftain? (\textbf{\kn{ಪಾಳೆಗಾರ}}, \emph{pāḷegāra}.) A mansion? (\textbf{\kn{ಮಹಲು}}, \emph{mahalu}.) To set foot? (\textbf{\kn{ಕಾಲಿಡು}}, \emph{kāliḍu}.) Exile in the forest? (\textbf{\kn{ವನವಾಸ}}, \emph{vanavāsa}.) To impose? (\textbf{\kn{ವಿಧಿಸು}}, \emph{vidhisu}.) Banishment? (\textbf{\kn{ಗಡೀಪಾರು}}, \emph{gaḍīpāru}.) Former? (\textbf{\kn{ಮಾಜಿ}}, \emph{māji}.) Theft? (\textbf{\kn{ಕಳ್ಳತನ}}, \emph{kaḷḷatana}.) A rise? (\textbf{\kn{ಉದಯ}}, \emph{udaya}.) A downfall? (\textbf{\kn{ಪತನ}}, \emph{patana}.) The past? (\textbf{\kn{ಗತಕಾಲ}}, \emph{gatakāla}.) A narrative? (\textbf{\kn{ವೃತ್ತಾಂತ}}, \emph{vṛttānta}.) A triumph? (\textbf{\kn{ವಿಜಯ}}, \emph{vijaya}.) Magic? (\textbf{\kn{ಜಾದೂ}}, \emph{jādū}.) A miracle? (\textbf{\kn{ಪವಾಡ}}, \emph{pavāḍa}.) False? (\textbf{\kn{ಹುಸಿ}}, \emph{husi}.) A noise? (\textbf{\kn{ಸದ್ದು}}, \emph{saddu}.) To kill? (\textbf{\kn{ಕೊಲ್ಲು}}, \emph{kollu}.) A camp? (\textbf{\kn{ಶಿಬಿರ}}, \emph{śibira}.) A refugee? (\textbf{\kn{ನಿರಾಶ್ರಿತ}}, \emph{nirāśrita}.) Condolence? (\textbf{\kn{ಸಂತಾಪ}}, \emph{santāpa}.) To hand over? (\textbf{\kn{ಹಸ್ತಾಂತರಿಸು}}, \emph{hastāntarisu}.) To rush? (\textbf{\kn{ಧಾವಿಸು}}, \emph{dhāvisu}.)
\end{tcolorbox}
\section[(dialogue)]{The Relief Camp — once more}

\label{lesson:KA-R375-review-final-375}



\begin{quote}
Say the Kannada for a camp and a refugee.
\end{quote}

\subsection*{Your turn}
\begin{itemize}
\raggedright
  \item \emph{Say it:} a camp, a refugee, condolence, to hand over, to rush
\end{itemize}

\begin{tcolorbox}[breakable,colback=teal!4,colframe=teal!35!black,title={Before you move on}]
A camp? (\textbf{\kn{ಶಿಬಿರ}}, \emph{śibira}.) A refugee? (\textbf{\kn{ನಿರಾಶ್ರಿತ}}, \emph{nirāśrita}.) Condolence? (\textbf{\kn{ಸಂತಾಪ}}, \emph{santāpa}.) To hand over? (\textbf{\kn{ಹಸ್ತಾಂತರಿಸು}}, \emph{hastāntarisu}.) To rush? (\textbf{\kn{ಧಾವಿಸು}}, \emph{dhāvisu}.)
\end{tcolorbox}
